\documentclass[11pt]{article}

\usepackage[margin=1in]{geometry}
\usepackage{amsmath}
\usepackage{amssymb}
\usepackage{amsthm}

\newtheorem{theorem}{Theorem}
\newtheorem{lemma}[theorem]{Lemma}
\newtheorem{proposition}[theorem]{Proposition}
\newtheorem{corollary}[theorem]{Corollary}
\theoremstyle{definition}
\newtheorem{conjecture}[theorem]{Conjecture}
\newtheorem{remark}[theorem]{Remark}

\newcommand{\R}{\mathbb{R}}
\newcommand{\Z}{\mathbb{Z}}
\newcommand{\T}{\mathbb{T}}
\newcommand{\Lam}{\Lambda}
\newcommand{\inner}[2]{\langle #1, #2\rangle}

\title{Refutation of Conjecture 00000000807:\\ the maximal multiplicity of $\lambda_1(\T^2)$ is $6$, not $4$}
\author{earthking11}
\date{13 September 2026}

\begin{document}
\maketitle

\begin{abstract}
We disprove Conjecture 00000000807.  For a flat two-torus $\R^2/\Lam$ the
Laplace eigenvalues are $4\pi^2|k|^2$ with $k$ ranging over the dual lattice
$\Lam^{*}$, so the multiplicity of the first non-zero eigenvalue $\lambda_1$
equals the number of shortest non-zero vectors of $\Lam^{*}$.  The square
torus has multiplicity $4$, but the triangular (equilateral) torus has
multiplicity $6$.  The true maximum over all flat tori is $6$, attained
exactly by lattices similar to the triangular lattice.  Consequently both
clauses of the conjecture fail.
\end{abstract}

\section{The conjecture}

The statement under consideration is:
\begin{quote}
\textbf{Conjecture 00000000807.} \emph{The maximal multiplicity of
$\lambda_1(\T^2)$ is $4$ (attained only on the square torus), and at most
$3$ for general lattices (multiplicity rigidity of lattice spectra,
partially known).}
\end{quote}
We show that this is false: the maximum is $6$, the square torus is not the
unique maximiser, and the bound ``at most $3$ for general lattices'' is
violated already by the square torus ($4$) and the triangular torus ($6$).

\section{Spectrum of a flat torus}\label{sec:spectrum}

Let $\Lam \subset \R^2$ be a full lattice and let
\[
  \Lam^{*} = \{\, k \in \R^2 : \inner{k}{\ell} \in \Z \ \text{for all } \ell \in \Lam \,\}
\]
be its dual lattice.  On the flat torus $\T^2 = \R^2/\Lam$ the Laplace
operator has the complete orthonormal system of eigenfunctions
\[
  \varphi_k(x) = e^{2\pi i \inner{k}{x}}, \qquad k \in \Lam^{*},
\]
with
\[
  -\Delta \varphi_k = 4\pi^2 |k|^2 \varphi_k .
\]
Hence the spectrum is $\{\,4\pi^2|k|^2 : k \in \Lam^{*}\,\}$, and for the
first non-zero eigenvalue
\begin{equation}\label{eq:mult}
  \operatorname{mult}(\lambda_1(\T^2))
  = \#\{\, k \in \Lam^{*} \setminus \{0\} : |k|^2 = r_1^2 \,\},
  \qquad
  r_1^2 = \min_{0 \neq k \in \Lam^{*}} |k|^2 .
\end{equation}
Since $k$ and $-k$ always have the same length, every multiplicity is even.

\section{The square torus has multiplicity $4$}

Take $\Lam = \Z^2$, so that $\Lam^{*} = \Z^2$ as well.  The shortest
non-zero vectors are
\[
  \pm e_1 = \pm(1,0), \qquad \pm e_2 = \pm(0,1),
\]
all of squared length $1$.  Thus
\[
  \operatorname{mult}(\lambda_1) = 4
\]
for the square torus, in agreement with the first numerical claim of the
conjecture.

\section{The triangular torus has multiplicity $6$}\label{sec:tri}

Let
\[
  \Lam = \Z(1,0) \oplus \Z\!\left(\tfrac12, \tfrac{\sqrt{3}}{2}\right)
\]
be the triangular (equilateral) lattice, with basis
$b_1 = (1,0)$, $b_2 = (\tfrac12, \tfrac{\sqrt3}{2})$.  Its Gram matrix is
\[
  G = \bigl(\inner{b_i}{b_j}\bigr)_{i,j}
    = \begin{pmatrix} 1 & 1/2 \\ 1/2 & 1 \end{pmatrix},
  \qquad \det G = \tfrac34 .
\]
The dual basis $d_1, d_2$, characterised by $\inner{d_i}{b_j} = \delta_{ij}$,
is
\[
  d_1 = \left(1, -\tfrac{1}{\sqrt3}\right), \qquad
  d_2 = \left(0, \tfrac{2}{\sqrt3}\right),
\]
and the Gram matrix of the dual lattice is
\[
  G^{*} = G^{-1}
        = \frac{4}{3}
          \begin{pmatrix} 1 & -1/2 \\ -1/2 & 1 \end{pmatrix}
        = \begin{pmatrix} 4/3 & -2/3 \\ -2/3 & 4/3 \end{pmatrix}.
\]
Explicitly, for $m,n \in \Z$,
\[
  |m d_1 + n d_2|^2
  = \frac{4}{3}\bigl(m^2 - mn + n^2\bigr).
\]
The shortest non-zero vectors $k \in \Lam^{*}$ are therefore the six vectors
\[
  \pm d_1, \qquad \pm d_2, \qquad \pm(d_1 + d_2),
\]
corresponding to $(m,n) = \pm(1,0), \pm(0,1), \pm(1,1)$, and each of them has
squared length
\[
  |d_1|^2 = |d_2|^2 = |d_1+d_2|^2 = \frac{4}{3}.
\]
The next possible length is strictly larger: for instance
$|d_1 - d_2|^2 = \frac43(1 + 1 + 1) = 4 > \frac43$.  By
\eqref{eq:mult},
\[
  \operatorname{mult}\bigl(\lambda_1(\T^2)\bigr) = 6
\]
for the triangular torus.  In particular the multiplicity is \emph{not}
bounded by $4$, and the maximum is \emph{not} attained only on the square
torus.

\begin{remark}
The triangular lattice is realised by the positive definite integral binary
quadratic form $Q(a,b) = a^2 + ab + b^2$: the six solutions of $Q = 1$ are
exactly $(1,0),(-1,0),(0,1),(0,-1),(1,-1),(-1,1)$.  This is the form whose
exact shortest-vector counts are formalised in \texttt{lean4/Main.lean}.
\end{remark}

\section{The true maximum is $6$}\label{sec:max}

The following standard argument (the two-dimensional kissing number, in the
form used for lattice spectra) shows that no flat torus has multiplicity
larger than $6$, and identifies the equality case.

\begin{theorem}\label{thm:max}
Let $\Lam \subset \R^2$ be any lattice and let
$r = r_1 = \min_{0 \neq k \in \Lam^{*}} |k|$.  Then the number of shortest
non-zero vectors of $\Lam^{*}$ is at most $6$.  If it equals $6$, then
$\Lam$ is similar to the triangular lattice.
\end{theorem}

\begin{proof}
Let $u \neq v$ be two distinct shortest non-zero vectors of $\Lam^{*}$, so
$|u| = |v| = r$.  Their difference $u - v$ is a non-zero vector of
$\Lam^{*}$, hence by minimality of $r$
\[
  |u-v|^2 \ge r^2 .
\]
Expanding,
\[
  |u-v|^2 = |u|^2 + |v|^2 - 2\inner{u}{v} = 2r^2 - 2\inner{u}{v} \ge r^2,
\]
so $\inner{u}{v} \le r^2/2$.  Writing $\inner{u}{v} = r^2 \cos\theta$ with
$\theta \in [0,\pi]$ the angle between $u$ and $v$, this says
$\cos\theta \le \tfrac12$, i.e.
\begin{equation}\label{eq:angle}
  \theta \ge 60^{\circ}.
\end{equation}
Thus the shortest vectors form $N$ points on the circle of radius $r$ whose
pairwise angular separations are at least $60^\circ$.  Order the points
cyclically around the circle and let $g_1,\dots,g_N$ be the $N$ successive
gaps, so that $\sum_i g_i = 360^\circ$.  If some gap were $< 60^\circ$, then
the two points bounding it would have angular separation $g_i < 60^\circ$
(their smaller angle is $g_i$, since $g_i < 180^\circ$), contradicting
\eqref{eq:angle}.  Hence every $g_i \ge 60^\circ$ and
$60^\circ N \le \sum_i g_i = 360^\circ$, that is $N \le 6$.

If $N = 6$, then all $N$ gaps equal $60^\circ$, so the six shortest vectors
are the vertices of a regular hexagon and the shortest-vector shell of
$\Lam^{*}$ is that hexagon.  This forces the lattice of shortest vectors to
be the hexagonal lattice, i.e. $\Lam$ is similar to the triangular lattice.
\end{proof}

\begin{corollary}
For a flat two-torus $\R^2/\Lam$ the multiplicity of $\lambda_1$ is one of
$2$, $4$, $6$.  It is $2$ for a generic lattice, $4$ for the square type
and $6$ for the triangular type; the maximum $6$ is attained exactly by
lattices similar to the triangular lattice.
\end{corollary}

\begin{proof}
Multiplicities are even, and by Theorem~\ref{thm:max} the shortest-vector
shell has at most $6$ vectors.  Since the shortest vectors come in $\pm$
pairs, the possible counts are $2$, $4$, $6$.  The square lattice realises
$4$ and the triangular lattice realises $6$; a generic lattice (with no
accidental coincidences between lengths) realises $2$.
\end{proof}

\section{Conclusion}

Both clauses of Conjecture 00000000807 are false.
\begin{enumerate}
  \item The maximal multiplicity of $\lambda_1(\T^2)$ is $6$, not $4$
        (Theorem~\ref{thm:max} and Section~\ref{sec:tri}).  The value $4$ is
        attained by the square torus, but it is not maximal.
  \item The square torus is not the unique lattice attaining the maximum:
        the triangular torus attains multiplicity $6 > 4$, and in fact the
        triangular torus is the unique maximiser up to similarity.
  \item The clause ``at most $3$ for general lattices'' is false as stated:
        the square torus has multiplicity $4$ and the triangular torus has
        multiplicity $6$.
\end{enumerate}

\section{Formal verification}

The exact shortest-vector counts are formalised in core Lean~4
(\texttt{lean4/Main.lean}, no Mathlib and no \texttt{sorry}):
\texttt{tri\_min} proves $a^2 + ab + b^2 \ge 1$ for every non-zero integer
pair; \texttt{tri\_six\_shortest} proves that $a^2 + ab + b^2 = 1$ has exactly
the six solutions listed above; \texttt{square\_four\_shortest} proves that
$a^2 + b^2 = 1$ has exactly four solutions; and
\texttt{conjecture\_00000000807\_false} records $4 < 6$ and the existence of a
multiplicity strictly larger than $4$.  The dual-lattice/spectrum
correspondence \eqref{eq:mult} and the $60^\circ$/kissing-number argument of
Theorem~\ref{thm:max} are proved here and reproduced numerically in
\texttt{reproduce.py}; they are not part of the Lean formalisation.

\end{document}
